\section{He Kaiming y FAIR: de ResNeXt a MAE}

Xie Saining (谢赛宁) describe su encuentro con He Kaiming (何凯明) como un punto de inflexión en su vida, y no exagera. La primera vez que colaboraron con intensidad, He Kaiming acababa de unirse a FAIR y todavía no sabía usar bien un Linux cluster, así que Xie Saining tuvo que enseñarle a trabajar en el clúster de cómputo. Este detalle resulta muy interesante, porque devuelve a He Kaiming —a quien después se consideraría un «investigador de nivel divino»— a una relación humana concreta: la gran investigación no surge de forma automática a partir de un genio abstracto, sino que crece en el proceso de resolver problemas codo a codo, ajustar experimentos y discutir cuestiones. Los dos participaron juntos en el ImageNet Challenge y terminaron produciendo ResNeXt, con el que obtuvieron el segundo lugar; esto no solo fue el punto de partida de un artículo clásico, sino también la primera vez que Xie Saining entró de verdad en la escena de «estar en sintonía con la investigación de primer nivel».

Él describe a He Kaiming como poseedor de un «campo gravitacional de distorsión de la realidad», es decir, que las personas a su alrededor son atraídas dentro de ese campo hacia un estándar más alto. Esta caracterización resulta precisa porque no se trata ni del charismatic leadership en el sentido tradicional, ni de un simple modelo de laboriosidad, sino de una capacidad para reforjar ideas ordinarias en «ideas de oro». Xie Saining dice: \textit{«Kaiming puede convertir todo lo muy ordinario en una idea de oro»}. Esto revela en realidad una capacidad central de los investigadores de primer nivel: no siempre están buscando elementos completamente novedosos; lo más común es que sean capaces de identificar, en material que otros consideran ordinario, una estructura que en verdad se puede amplificar.

\begin{importantbox}{ResNeXt: cómo un componente ordinario se depuró hasta convertirse en un paradigma}
De ResNeXt suele recordarse la convolución agrupada y una mejor relación entre precisión y eficiencia, pero el punto que Xie Saining subraya es que su idea es, en esencia, cercana al MoE actual: no se trata de aumentar sin más la profundidad o la anchura, sino de introducir en la estructura la idea de «varios expertos en paralelo que después se agregan». Por eso, la «X» tiene tanto el sentido de next como, de manera velada, el nombre de Xie escondido en ella, como una especie de recordatorio de una creación compartida.
\end{importantbox}

El verdadero impacto de ResNeXt en Xie Saining no fue solo el éxito a nivel de artículo, sino que le permitió ver cómo investigaba He Kaiming: partir de una variante estructural que no era llamativa, captar el grado de libertad más esencial y después pulirlo hasta convertirlo en un diseño a la vez simple y con poder explicativo. Este estilo se extendió más adelante a lo largo del trabajo de aprendizaje autosupervisado del periodo en FAIR. Yann LeCun ya había comparado el self-supervised learning con la base de un pastel, al considerar que la mayor parte de la inteligencia proviene de predecir y representar la estructura del mundo, y no de una supervisión realizada con pocas etiquetas. El MoCo y el MAE que Xie Saining y He Kaiming impulsaron en FAIR fueron precisamente intentos que, en esta dirección general, se acercaban cada vez más a algo «más simple, más scalable y más cercano a lo esencial».

Hoy en día es fácil subestimar el valor de MoCo, porque conceptos como el aprendizaje contrastivo, el preentrenamiento autosupervisado y el codificador de momento se volvieron después algo muy conocido. Pero, en su momento, MoCo era importante porque, por primera vez, hizo que el aprendizaje contrastivo work de manera estable en el preentrenamiento visual. Estableció un mecanismo que se actualizaba de manera continua y que podía aprovechar muestras negativas a gran escala, con lo cual llevó los métodos de aprendizaje de representaciones, que antes eran más frágiles, a una etapa más práctica. El MAE, en cambio, representó otro giro: en lugar de seguir haciendo más complejo el aprendizaje contrastivo, volvió a la forma simple del masked prediction y descubrió que «menos supuestos» tenía, por el contrario, más potencial en el escalamiento. Este cambio de MoCo a MAE reflejaba, en sí mismo, un temperamento de investigación propio de FAIR: no desarrollar un apego emocional a la línea de trabajo ya establecida y estar dispuestos a barajar todo de nuevo ante evidencia nueva.

\begin{knowledgebox}{La continuidad entre MoCo y MAE}
MoCo y MAE parecen métodos muy distintos, pero comparten la misma conciencia del problema: cómo lograr que un sistema de visión aprenda, a partir de enormes cantidades de datos, estructuras transferibles, sin depender de etiquetas hechas por humanos. MoCo se parece más a «aprender representaciones mediante alineación discriminativa», mientras que MAE se parece más a «hacer emerger la semántica interna mediante la reconstrucción de lo faltante». Ninguno de los dos es un punto final, pero juntos demuestran que el aprendizaje autosupervisado no es una rama secundaria en la visión, sino la línea principal.
\end{knowledgebox}

Sin embargo, fue precisamente en el periodo en FAIR cuando Xie Saining fue tocando poco a poco un hecho más difícil: que el aprendizaje autosupervisado «funcione» no equivale automáticamente a que «se pueda escalar [scale] de manera ilimitada como los modelos de lenguaje». Menciona que He Kaiming ya decía muy pronto, entre 2018 y 2019, que «había que hacer los modelos cada vez más y más grandes», pero, conforme acumularon experiencia, descubrieron que la autosupervisión visual funcionaba en todos los domain, aunque era muy difícil que apareciera una scaling law tan suave, sostenida y deslumbrante como la de los LLM. Este hallazgo es importante porque obliga a los investigadores a reconocer que la estructura de supervisión, la densidad de información y la distribución de ruido del mundo visual no son simétricas respecto al texto. Precisamente por eso, es muy probable que los futuros modelos del mundo no repitan sin más la trayectoria de crecimiento de los modelos de lenguaje.

\begin{warningbox}{«Hacerlo más grande siempre mejora» no es una verdad que cruce modalidades}
Uno de los contraejemplos más memorables de la época de FAIR es que la autosupervisión visual no reprodujo de manera natural la curva del ámbito del lenguaje solo porque los modelos fueran más grandes y hubiera más datos. La experiencia de Xie Saining nos recuerda que el scaling no es una llave universal, sino un fenómeno que depende necesariamente de la estructura del problema. Ignorar las diferencias entre modalidades solo lleva a proyectar de manera equivocada el mito del éxito de los modelos de lenguaje sobre todas las direcciones.
\end{warningbox}

Hay otro nivel de continuidad que a menudo se pasa por alto: el trabajo de autosupervisión del periodo en FAIR no «quedó obsoleto» tras la aparición de DiT o de los modelos del mundo, sino que preparó de antemano la metodología y la intuición para todos los giros posteriores. MoCo y MAE entrenaron a los investigadores para creer que las etiquetas no son la única señal de aprendizaje, y que lo verdaderamente valioso es cómo construir la tarea de modo que el modelo se vea obligado a comprimir estructuras estables del mundo. Cuando Xie Saining, más adelante, empezó a pensar en el predictive world model, esta convicción se extendió de manera natural hacia allá. Dicho de otro modo, él no saltó de golpe del «aprendizaje de representaciones» al «modelo del mundo», sino que siguió avanzando río abajo dentro de la misma corriente de pensamiento.

El significado de trabajar cerca de He Kaiming está también en que fue la primera vez que Xie Saining sintió de manera sistemática cómo se transmite un «estilo de investigación». He Kaiming no enseñaba la metodología mediante adoctrinamiento verbal, sino que, a través de experimento tras experimento y detalle de escritura tras detalle de escritura, convertía la indagación sobre la esencia del problema en el estándar de trabajo predeterminado del equipo. Que Xie Saining, más adelante, pudiera hablar con tanta claridad de research taste, de la exploración no lineal y de la narrativa de los artículos se debe, en buena medida, a que en FAIR presenció una máquina de investigación de alto nivel funcionando de manera integral.

\subsection{Resumen de la sección}

De ResNeXt a MoCo y MAE, lo que Xie Saining vivió en FAIR no fue el encadenamiento de unos cuantos artículos representativos, sino un remodelado de su metodología de investigación. He Kaiming le mostró que el trabajo de primer nivel suele surgir de recombinar elementos ordinarios, de la simplificación continua de la línea principal del aprendizaje autosupervisado y de corregir a tiempo las ilusiones sobre el scaling.
